% Lección: álgebra de límites (castellano).
%
% La demostración casi nunca cabe en una diapositiva, así que va en
% `notesonly`: sale en los apuntes y desaparece de las diapositivas sin que
% haga falta un segundo fichero.

\begin{frame}
\didactatitle{Álgebra de límites}

\begin{proposition}[Álgebra de límites]
Si $\lim_{x \to a} f(x) = L$ y $\lim_{x \to a} g(x) = M$, entonces
\begin{itemize}
  \item $\lim_{x \to a} \bigl(f(x) + g(x)\bigr) = L + M$;
  \item $\lim_{x \to a} f(x)\,g(x) = L M$;
  \item $\lim_{x \to a} \dfrac{f(x)}{g(x)} = \dfrac{L}{M}$, si $M \neq 0$.
\end{itemize}
\end{proposition}

\begin{notesonly}
\begin{proof}
Probamos la primera; las otras dos siguen la misma idea con alguna acotación
más. Dado $\varepsilon > 0$, existen $\delta_1, \delta_2 > 0$ tales que
$|f(x) - L| < \varepsilon/2$ si $0 < |x - a| < \delta_1$, y
$|g(x) - M| < \varepsilon/2$ si $0 < |x - a| < \delta_2$. Tomando
$\delta = \min\{\delta_1, \delta_2\}$, para $0 < |x - a| < \delta$ la
desigualdad triangular da
\[
  \bigl|\bigl(f(x) + g(x)\bigr) - (L + M)\bigr|
  \le |f(x) - L| + |g(x) - M|
  < \tfrac{\varepsilon}{2} + \tfrac{\varepsilon}{2} = \varepsilon .
\]
\end{proof}
\end{notesonly}

\dpause

\begin{remark}
Como $\lim_{x \to a} x = a$ y el límite de una constante es ella misma, la
proposición da el límite de un polinomio evaluando: $\lim_{x \to a} p(x) =
p(a)$. Y el de una función racional, $\lim_{x \to a} p(x)/q(x) = p(a)/q(a)$,
siempre que $q(a) \neq 0$.
\end{remark}

\timing{15 min}
\end{frame}
